\documentclass[12pt]{article}
\usepackage{graphicx}
\usepackage{amsmath}
\usepackage{hyperref}
\usepackage{natbib}

\title{The Impact of Urbanization on Biodiversity: Habitat Loss, Species Extinction, and Ecosystem Services in Urban Environments}
\author{Your Name}
\date{\today}

\begin{document}

\maketitle

\begin{abstract}
Urbanization significantly impacts biodiversity through habitat loss, species extinction, and the alteration of ecosystem services. This paper analyzes these impacts from a global perspective, reviews case studies from different parts of the world, and proposes mitigation strategies supported by extensive literature and data. Recommendations for future urban planning policies to minimize biodiversity loss are also provided.
\end{abstract}

\section{Introduction}
Urbanization, the process by which rural areas are transformed into urban areas, is a key driver of environmental change. It results in significant alterations to land use and has profound effects on biodiversity. This paper examines the impacts of urbanization on biodiversity, focusing on habitat loss, species extinction, and ecosystem services in urban environments. A comprehensive review of pertinent literature and case studies is provided, followed by proposed strategies to mitigate negative impacts.

\section{Impact Analysis}

\subsection{Habitat Loss}
One of the primary impacts of urbanization on biodiversity is habitat loss. Land conversion for urban development often leads to the destruction of natural habitats. According to \citet{seto2012}, urban expansion is expected to threaten over 30,000 km\textsuperscript{2} of biodiversity hotspots by 2030.

\subsection{Species Extinction}
As habitats are lost and fragmented, species that cannot adapt to urban environments face heightened risks of extinction. Urbanization has been linked to the decline of numerous species worldwide. \citet{mcdonald2008} estimated that urbanization could lead to the extinction of approximately 8\% of vertebrate species.

\subsection{Ecosystem Services}
Urban environments also affect ecosystem services, which are critical for human well-being. Key services such as air and water purification, climate regulation, and pollination are disrupted by urban sprawl. \citet{bolund1999} discuss how urban green spaces can mitigate some of these impacts by providing essential ecosystem services within cities.

\section{Case Studies}

\subsection{North America}
In North America, cities like New York and Los Angeles show significant biodiversity changes over the last century. Habitat fragmentation and pollution have led to the decline of species such as the monarch butterfly (\textit{Danaus plexippus}).

\subsection{Asia}
Asian megacities like Tokyo and Beijing have experienced rapid urbanization, which has led to the loss of forested areas and wetlands. This has resulted in declines in amphibian and bird species as discussed by \citet{saito2021}.

\subsection{Africa}
In Africa, cities such as Nairobi have seen encroachment into wildlife corridors and national parks that provide critical habitats for large mammals including elephants and lions. According to \citet{okanga2017}, urbanization in these regions threatens the last remaining habitats of several endangered species.

\section{Mitigation Strategies}
To mitigate the negative impacts of urbanization on biodiversity, several strategies can be implemented:

\subsection{Urban Green Spaces}
The creation and maintenance of urban green spaces can help preserve local biodiversity within cities. Green roofs, parks, and community gardens provide habitats for urban wildlife and contribute to ecosystem services.

\subsection{Urban Planning}
Sustainable urban planning practices, such as promoting high-density development and protecting natural areas, can reduce the footprint of urban expansion. Integrative approaches that balance development with conservation goals are essential.

\subsection{Legislation and Policy}
Strong environmental policies and regulations are crucial. Implementing zoning laws that restrict development in ecologically sensitive areas and requiring environmental impact assessments for urban projects can prevent habitat destruction.

\section{Recommendations}
Future urban governance should incorporate biodiversity considerations into planning processes. Multidisciplinary collaborations between ecologists, urban planners, and policymakers are necessary. Furthermore, public awareness and engagement in conservation efforts can enhance local biodiversity conservation initiatives.

\section{Conclusion}
Urbanization presents significant challenges to biodiversity, but with careful planning and management, it is possible to mitigate its adverse effects. This paper highlights the importance of preserving habitats, preventing species extinction, and maintaining ecosystem services in urban environments. Implementing the proposed strategies and recommendations will be essential for sustainable urban development that respects and protects biodiversity.

\bibliographystyle{apalike}
\bibliography{prompt13_4o}

\end{document}